% Generated by revision/scripts/loto_table.py
\begin{table}[!h]
\centering\small
\caption{Leave-one-tissue-out. Each row holds out one tissue
entirely. The across-tissue standard deviation is the quantity a
tissue-blind error bar should convey, and it is several times the
seed-to-seed standard deviation measured on any one fold, so a
single tissue carrying a seed-derived error bar understates how
far the number moves.}
\label{tab:loto}
\begin{tabular*}{\columnwidth}{@{\extracolsep{\fill}}l r rr@{}}
\toprule
\textbf{Held-out tissue} & \textbf{Test rows} & \textbf{PathXDRP} & \textbf{CDRScan} \\
\midrule
lung adeno. (LUAD) & $6{,}026$ & $0.8562$ & $0.8733$ \\
breast & $4{,}883$ & $0.8082$ & $0.8421$ \\
large intestine & $4{,}606$ & $0.8293$ & $0.8848$ \\
small-cell lung & $4{,}176$ & $0.8605$ & $0.8744$ \\
ovary & $3{,}793$ & $0.8601$ & $0.8876$ \\
\midrule
mean across tissues & & $0.8428$ & $0.8725$ \\
\textbf{SD across tissues} & & $0.0233$ & $0.0181$ \\
SD across seeds, fold 0 & & $0.0045$ & $0.0019$ \\
\bottomrule
\end{tabular*}

\smallskip
\footnotesize\emph{Note:} One seed per cell. The across-tissue
spread is a property of which tissue is held out and is not reduced
by averaging over seeds.
\end{table}